
\begin{proposition}[CDP20 Proposition 7.3 with a complete replacement proof]\label{cdp:st:A.4.9}
Let $Y$ be a connected smooth compact complex threefold with $H^1(Y,\Z)=H^2(Y,\Z)=0$, algebraic dimension one and holomorphic algebraic reduction $g:Y\to C$. There is a dense open subset of $\Pic(Y)$, whose complement is a proper complex analytic subset, such that
\[
 H^0\bigl((Y_c)_{\rm red},
       \widetilde\Omega^1_{(Y_c)_{\rm red}}\otimes L|_{(Y_c)_{\rm red}}\bigr)=0
 \tag{A.23}\label{cdp:eq:A.23}
\]
for every singular reduced fibre support and every $L$ in that subset. In fact the transgression twists $L_r$ of \ref{cdp:st:D.1.9} satisfy \eqref{cdp:eq:A.23} on every reduced fibre support, including smooth and reducible supports, for every real irrational $r$.
\end{proposition}
\begin{proof}
The independently rederived smooth-fibre classification \ref{cdp:st:C.3.1} leaves complex tori, bielliptic surfaces, primary or secondary Kodaira surfaces, and primary or secondary Hopf surfaces or Inoue surfaces. We prove that for each such smooth fibre $F$ and each nonzero $\gamma_F\in H^1(F,\Z)$ the local system
\[
 \mathbb A_{r,F}:\quad a\longmapsto
       \exp(2\pi\mathrm i r\langle\gamma_F,a\rangle)
 \quad(r\in\mathbb R\setminus\Q)
\]
satisfies $H^0(F,\mathbb A_{r,F})=H^1(F,\mathbb A_{r,F})=0$. Its character has infinite image, so the assertion in degree zero is immediate.

For a torus the Koszul argument in \ref{cdp:st:A.4.8} applies. A bielliptic surface has a finite étale torus cover. Pullback of the integral class is nonzero: transfer composed with pullback is multiplication by the degree. The character stays nontrivial after pullback, and transfer with local coefficients makes pullback on cohomology injective over $\C$. Thus degree-one vanishing descends from the torus cover.

A primary Kodaira surface is a principal elliptic bundle over an elliptic curve, as described in \cite[\S1 and the central-extension diagram in the proof of Lemma~1.1]{SS2002}. Its fundamental group has a central extension
\[
 1\longrightarrow\Z^2\longrightarrow\Gamma\longrightarrow\Z^2
       \longrightarrow1.
 \tag{A.24}\label{cdp:eq:A.24}
\]
For any nontrivial one-dimensional character $\chi$ of $\Gamma$, if its restriction to the central subgroup is nontrivial, the latter's cohomology with these coefficients is zero by the two-variable Koszul contraction. The Hochschild--Serre spectral sequence then gives $H^1(\Gamma,\C_\chi)=0$. If its central restriction is trivial, $\chi$ descends to a nontrivial character of the quotient $\Z^2$. Centrality makes its action on the cohomology of the central subgroup equal to that scalar character, and the same Koszul contraction kills every term on the second page. Again $H^1(\Gamma,\C_\chi)=0$. First cohomology of a connected space with local coefficients is first cohomology of its fundamental group, so this proves the required assertion. A secondary Kodaira surface has a finite primary cover; the same transfer argument as above applies. Hopf surfaces have finite primary covers, and a primary Hopf surface is diffeomorphic to $S^1\times S^3$; the cover statement and topology are recorded in \cite[\S2.1.1 and Theorem~2.2]{Albanese2020}. The nontrivial circle character has an invertible cellular differential, giving the required vanishing and its finite-cover descent.

For completeness we also give the Inoue calculation, rather than assuming a generic-vanishing theorem for its cotangent bundle. The explicit constructions in \cite[\S2.1.2 and Proposition~2.3]{Albanese2020} identify every Inoue surface with a mapping torus $N_\varphi$ of a compact three-manifold $N$. In the $S_M$ cases $N$ is a three-torus and the eigenvalues of $\varphi^*$ on $H^1(N,\C)$ are $\alpha,\beta,\bar\beta$, with $\alpha>1$ and $|\beta|=\alpha^{-1/2}<1$, after reversing the circle orientation when necessary. In the $S^+$ and $S^-$ cases $N$ is the nontrivial circle bundle over a two-torus; the nonzero Euler class makes $H^1(N,\C)$ the pullback of $H^1(T^2,\C)$ by its Gysin sequence. The monodromy on it has eigenvalues $\alpha,\alpha^{-1}$ or $\alpha,-\alpha^{-1}$, respectively. These assertions follow as well directly from the conjugation relations displayed in that construction: on the abelian rank-three kernel the matrix is $M$, and on the quotient of the Heisenberg kernel by its centre the matrix is $N\in\mathrm{GL}_2(\Z)$ with precisely those eigenvalues.

None of these degree-one eigenvalues has modulus one. In particular there are no invariant degree-one classes on $N$. The ordinary mapping-torus sequence therefore identifies $H^1(F,\Z)$ with the copy of $\Z$ from the circle, since the invariant subgroup of the torsion-free group $H^1(N,\Z)$ is zero. Thus $\mathbb A_{r,F}$ is pulled back from the circle, with monodromy $\lambda\in S^1$ of infinite order. The local-coefficient Leray sequence for $N\to N_\varphi\to S^1$ gives
\[
\begin{split}
0\longrightarrow
 \coker(\lambda-1:H^0(N,\C)\to H^0(N,\C))
 \longrightarrow H^1(F,\mathbb A_{r,F})\\
 \longrightarrow
 \ker(\lambda\varphi^*-1:H^1(N,\C)\to H^1(N,\C))
 \longrightarrow0.
\end{split}
\]
The first term is zero because $\lambda\ne1$, and the last term is zero by the eigenvalue moduli. Using the reciprocal convention for the circle local system replaces $\lambda$ by $\lambda^{-1}$ and gives the same conclusion. This completes the degree-one calculation for every surface type in \ref{cdp:st:C.3.1}.

Now repeat the neighbourhood argument of \ref{cdp:st:A.4.8} with these two fibre-cohomology vanishings. \ref{cdp:st:D.1.8}--\ref{cdp:st:D.1.9} supply the same nonzero integral class on every nearby smooth fibre and the unitary local system on every fibre neighbourhood. The only terms that can contribute to first cohomology of the punctured neighbourhood are its degree-zero and degree-one fibre cohomology, just proved zero. The fundamental-group surjection and proper base change therefore give $H^1((Y_c)_{\rm red},\mathbb A_{r,c})=0$ for every $c$. Applying \ref{cdp:st:A.4.7} proves \eqref{cdp:eq:A.23} for the single global twist $L_r$, simultaneously on all reduced supports.

Finally, the exponential sequence and the integral vanishings identify $\Pic(Y)$ with the complex vector space $H^1(Y,\OO_Y)$. Choose holomorphic Cech cocycles representing a basis; exponentiating their linear combinations gives a holomorphic family of line bundles on $Y\times H^1(Y,\OO_Y)$ whose fibre is the indicated line bundle. There are finitely many singular fibre supports. For each of them, tensor its fixed coherent torsion-free cotangent sheaf by this family and apply proper semicontinuity: the locus where its space of sections is nonzero is a closed complex analytic subset of the parameter space. Each is proper because it excludes $L_r$. Their finite union is proper and has empty interior, giving the asserted dense open set. This proves the statement of Proposition~7.3. The invalid genus-two deduction in its printed proof remains refuted by \ref{cdp:st:A.4.2}; the proof above replaces that argument through the explicit maps (A.20)--\eqref{cdp:eq:A.24}.
\end{proof}


